\chapter{Sugar}
\index{Sugar}

\medskip
\noindent\textit{Educational reference; always seek professional advice for individual care.}

\section*{Definition and Overview}
Sugar is the common name for a group of sweet-tasting, soluble carbohydrates. In nutrition, the word is used in two ways: for refined white and brown sugar, honey, syrups, and the sugars added to foods during processing, and for the natural sugars that occur in fruit, vegetables, and milk. Sugars are an energy source for the body, and every cell uses glucose, the simplest sugar, for fuel. The issue is not that sugar exists; it is how much, in what form, and how often we consume it. The body handles the sugars in whole fruit and vegetables very differently from the free sugars added to soft drinks, biscuits, cakes, and processed foods, and it is the free sugars that cause the most harm.

In many countries, as in most developed countries, people eat far more free sugar than is recommended, and most of it is hidden in processed foods rather than added at the table. Excessive sugar intake is linked to obesity, type 2 diabetes, tooth decay, and heart disease, making it one of the most important dietary issues of modern public health. The relationship between sugar and health is not a matter of a single food being "good" or "bad"; it is a question of quantity, frequency, and the form in which sugar is consumed. This chapter explains what sugars are, how the body processes them, what the local recommendations say, and where sugar hides in the food supply. It ends with practical, realistic advice for reducing intake without turning eating into a chore. The aim is not perfection or a sugar-free life, but a healthier balance that most families can actually maintain, and the evidence shows that even modest reductions in added sugar produce meaningful improvements in weight, dental health, and long-term disease risk.

\section*{What Counts as Sugar}
\begin{itemize}
    \item \textbf{Free sugars:} sugars added to foods and drinks by the manufacturer, cook, or consumer, plus the sugars naturally present in honey, syrups, fruit juice, and juice concentrates; these are the ones health bodies advise limiting
    \item \textbf{Naturally occurring sugars:} sugars that are built into the structure of whole fruit, vegetables, and dairy foods; these are not counted against the free-sugar limit because they come packaged with fibre, water, and nutrients
    \item \textbf{Glucose and fructose:} the two simplest sugars; glucose is the body's main fuel, and fructose is the sugar that makes fruit taste sweet
    \item \textbf{Sucrose:} common table sugar, a molecule made of one glucose and one fructose linked together
    \item \textbf{Lactose:} the sugar found naturally in milk
    \item \textbf{Maltose:} the sugar formed when starch is digested, found in malted products
    \item \textbf{Sugar alcohols:} sweeteners such as sorbitol and xylitol, which provide sweetness with fewer calories but can cause bloating and diarrhoea in large amounts
\end{itemize}

\section*{How the Body Handles Sugar}
When we eat sugar, the body breaks it down into glucose, which is absorbed into the bloodstream and carried to the cells, where it is converted into energy. The hormone insulin, produced by the pancreas, acts as the key that allows glucose into the cells and keeps blood sugar levels in a healthy range. Foods differ in how quickly they raise blood sugar: a sugary drink spikes it fast, while an apple, with its fibre, releases it slowly. Frequent large spikes force the pancreas to work harder, and over time the body's cells can become resistant to insulin, a precursor to type 2 diabetes.

Sugar that is not immediately used for energy is stored, first in the liver and muscles as glycogen and then, in larger amounts, as fat. This is why excess sugar, particularly the liquid sugar in soft drinks, is strongly associated with weight gain. Crucially, sugary drinks do not satisfy appetite in the way solid food does, so people tend to consume them in addition to, rather than instead of, their normal food. Fructose is processed almost entirely by the liver, and a high intake of fructose, especially in liquid form, contributes to fatty liver and raised blood fats. The body has no nutritional requirement for added sugar at all.

\section*{Health Effects of Too Much Sugar}
\begin{itemize}
    \item \textbf{Weight gain and obesity:} excess sugar adds calories without fullness, and the rising consumption of sugary drinks is a major driver of the obesity epidemic
    \item \textbf{Type 2 diabetes:} chronic excess sugar and weight gain increase insulin resistance, and people who drink sugar-sweetened beverages regularly are at substantially higher risk
    \item \textbf{Tooth decay:} the bacteria in dental plaque feed on sugar and produce acid that dissolves tooth enamel; every exposure to sugar, especially between meals, raises the risk of cavities
    \item \textbf{Heart disease:} high intakes of added sugar raise blood pressure, blood fats, and inflammation, all of which increase cardiovascular risk
    \item \textbf{Fatty liver:} excess fructose is converted to fat in the liver, leading to non-alcoholic fatty liver disease even in people who are not overweight
    \item \textbf{Gout:} fructose raises uric acid levels, and sugary drinks are associated with attacks of gout
    \item \textbf{Energy and mood:} rapid rises and falls in blood sugar can cause tiredness, irritability, and cravings
    \item \textbf{Displacement of nutritious food:} sugary foods crowd out the fruit, vegetables, and proteins that provide genuine nourishment
\end{itemize}

\section*{How Much Sugar Is Recommended}
local and World Health Organization guidance is clear about free sugar: it should be limited to less than 10 percent of daily energy intake, and ideally below 5 percent for additional health benefits. For an average adult, 10 percent is roughly 50 grams of free sugar a day, about 12 teaspoons, and 5 percent is about 25 grams, or 6 teaspoons. To put that in perspective, a single 600 ml bottle of soft drink typically contains about 16 teaspoons of sugar, which exceeds the entire daily allowance on its own. Children should aim for the same proportions of a smaller intake, and the advice is to avoid giving free sugar to children under two years at all.

The local Dietary Guidelines translate this into practical advice: limit intake of foods and drinks containing added sugars, drink plenty of water, and choose water rather than sugary drinks. Most people exceed the limit, largely through soft drinks, fruit juice, sports drinks, biscuits, cakes, breakfast cereals, flavoured yoghurts, and confectionery. National surveys show that sugary drinks are consumed by a large share of children and adults each week, and that the average daily intake of free sugar sits well above the recommended ceiling. Intake tends to be highest among teenagers and young adults, and the amount of sugar consumed in liquid form is a particular concern, because liquid sugar is absorbed quickly, does not register as food in the appetite system, and is strongly linked to weight gain. The recommendation is not about a single treat; it is about the daily pattern, because the total amount and the frequency of exposure determine the harm.

\section*{Where Sugar Hides}
\begin{itemize}
    \item \textbf{Sugary drinks:} soft drinks, fruit juice, cordial, energy drinks, and sports drinks; the single largest source of free sugar in the local diet
    \item \textbf{Breakfast cereals:} many cereals marketed to children contain several teaspoons of sugar per bowl
    \item \textbf{Flavoured yoghurts and dairy desserts:} can contain as much sugar as soft drink
    \item \textbf{Baked goods and confectionery:} biscuits, cakes, muffins, pastries, lollies, and chocolate
    \item \textbf{Sauces and spreads:} tomato sauce, barbecue sauce, sweet chilli, and jam add sugar to otherwise savoury meals
    \item \textbf{Seemingly healthy foods:} muesli bars, protein bars, dried fruit snacks, and flavoured oats are often dense with added sugar
    \item \textbf{Alcohol:} sweet wines, premixed drinks, and cider contain substantial sugar
    \item \textbf{Hidden names:} sugar appears on labels as sucrose, glucose, fructose, honey, molasses, agave, maple syrup, corn syrup, malt, and dozens of other terms ending in "ose"
\end{itemize}

\section*{When to Seek Professional Advice}
Most people can reduce sugar intake with the practical steps below, but professional advice is worthwhile in several situations. A general practitioner or an accredited practising dietitian can help if you are struggling with weight, if you have been told your blood sugar is high, or if you have prediabetes or diabetes. A dietitian can also help with a child's diet, with healthy eating during pregnancy, and with special situations such as diabetes in pregnancy or gestational diabetes.

Seek medical attention promptly if you develop symptoms that might suggest diabetes or another health problem, such as:
\begin{itemize}
    \item Increased thirst and urination, unexplained weight loss, or constant tiredness
    \item Blurred vision or slow healing of cuts and infections
    \item Recurrent boils or skin infections
    \item Recurrent tooth decay, which a dentist should assess
\end{itemize}

\textbf{Contact your local emergency services for:} severe drowsiness, confusion, rapid breathing, or collapse, which can occur with dangerously high or low blood sugar, especially in someone known to have diabetes.

\section*{Risks and Cautions}
\begin{itemize}
    \item \textbf{Diabetic ketoacidosis:} in people with diabetes, very high blood sugar can cause a life-threatening emergency requiring urgent hospital care
    \item \textbf{Sudden withdrawal effects:} cutting sugar abruptly can cause temporary headaches, tiredness, and cravings, which usually settle within a week or two
    \item \textbf{Low-calorie sweeteners:} artificial sweeteners are a reasonable tool for some people, but reliance on them does not necessarily change food habits, and large amounts of sugar alcohols can cause bloating and diarrhoea
    \item \textbf{Extreme restriction:} rigid "no sugar" rules can create anxiety and an unhealthy relationship with food, and can backfire into binge eating
    \item \textbf{Unreliable health claims:} products labelled "natural", "raw", or "unrefined" sugar are still sugar and count towards the daily limit
    \item \textbf{Fruit juice is not the same as fruit:} even unsweetened juice counts as free sugar because the fibre has been removed, and whole fruit is the healthier choice
\end{itemize}

\section*{Tips and Practical Advice}
\begin{itemize}
    \item Make water the default drink, and serve it with meals; add fruit slices or cucumber for flavour
    \item Choose whole fruit over juice and over dried fruit, which is concentrated sugar
    \item Read nutrition labels and check the "sugars" line, remembering that the free-sugar limit is around 25 to 50 grams a day
    \item When choosing packaged foods, compare products and pick the one with less added sugar
    \item Reduce sugar in recipes gradually, cutting the amount in baking and cooking by a quarter to a third over time
    \item Keep treats as occasional foods rather than everyday ones, and have healthy snacks such as nuts, fruit, and plain yoghurt on hand
    \item Start small: one achievable change, such as swapping soft drink for sparkling water, is more sustainable than a complete overhaul
    \item Involve the whole family, because children learn eating habits from the adults around them
    \item If you have diabetes or high blood sugar, follow your doctor's and dietitian's plan rather than experimenting with diets alone
    \item Be kind to yourself: occasional treats are part of a normal, enjoyable diet, and perfection is neither required nor helpful
\end{itemize}

\section*{Sources and Further Reading}
The following organisations provide reliable information about sugar, healthy eating, and the Australian dietary guidelines.

\begin{itemize}
    \item Australian Government Department of Health and Aged Care. \textit{Australian Dietary Guidelines and advice on sugar.} https://www.health.gov.au/
    \item healthdirect Australia. \textit{Sugar and healthy eating information.} https://www.healthdirect.gov.au/
    \item World Health Organization (WHO). \textit{Sugar intake for adults and children: guideline.} https://www.who.int/
    \item NHS. \textit{How to cut down on sugar in your diet.} https://www.nhs.uk/
    \item Mayo Clinic. \textit{Added sugar: where it hides and how to limit it.} https://www.mayoclinic.org/
    \item Dietitians Australia. \textit{Healthy eating and sugar information.} https://dietitiansaustralia.org.au/
\end{itemize}

\clearpage
